\documentclass[12pt]{article}
% Préambule commun à tous les documents extraits de « La Revue CAPES »
\usepackage[T1]{fontenc}
\usepackage[utf8]{inputenc}
\usepackage{lmodern}
\usepackage{amsmath,amssymb,amsthm,amsfonts,mathrsfs}
\usepackage{setspace}
\usepackage{listings}
\usepackage{pgfplots}
\pgfplotsset{compat=1.18}
\usepackage{longtable}
\usepackage{adjustbox}
\usepackage{lettrine}
\usepackage{tkz-tab}
\usepackage{changepage}
\usepackage{pdflscape}
\usepackage{graphicx}
\usepackage{tikz}
\usepackage{xcolor}
\usepackage[a4paper,top=2cm,bottom=2.5cm,left=2.5cm,right=2cm]{geometry}
\usepackage{fancyhdr}
\usepackage{draftwatermark}
\SetWatermarkText{La RevueCapes}
\SetWatermarkLightness{0.9}
\SetWatermarkScale{2}
\usepackage{hyperref}
\hypersetup{colorlinks=true,linkcolor=blue!50!black,urlcolor=blue!50!black}

\definecolor{revue}{RGB}{20,60,120}

\pagestyle{fancy}
\fancyhf{}
\renewcommand{\headrulewidth}{0.4pt}
\setlength{\headheight}{15pt}

% \entete{Matière}{Titre}{Type de document}
\newcommand{\entete}[3]{%
  \fancyhead[L]{\small\textsc{La Revue CAPES} -- #1}%
  \fancyhead[R]{\small #2}%
  \fancyfoot[C]{\thepage}%
  \thispagestyle{plain}%
  \begin{center}
    {\color{revue}\rule{\textwidth}{1.2pt}}\\[4pt]
    {\small\textsc{Concours directs CAP-CEG / CAPES -- Burkina Faso}}\\[6pt]
    {\LARGE\bfseries\color{revue} #2}\\[6pt]
    {\large\itshape #1 -- #3}\\[2pt]
    {\color{revue}\rule{\textwidth}{1.2pt}}
  \end{center}
  \vspace{0.5cm}%
}

% Encadré pour les remarques ajoutées dans les corrigés
\newcommand{\remarque}[1]{\par\medskip\noindent\fcolorbox{revue}{revue!6}{\parbox{\dimexpr\linewidth-2\fboxsep-2\fboxrule}{\textbf{\color{revue}Remarque.} #1}}\par\medskip}


\begin{document}
\entete{Mathématiques}{Sujet type n°12}{Énoncé}
\begin{spacing}{1.5}

\medskip{\textbf{Exercice 1}

1. Calculer le volume du domaine : 

$D=\lbrace (x,y,z)\in \mathbb{R}^3;0\leq x\leq 1;0\leq y\leq (1-\sqrt{x})^2;z^2\leq 2xy\rbrace$.

2. Soit $D=\lbrace (x,y)\in \mathbb{R}^2;x+y\leq 1;x\geq 0;y\geq 0\rbrace$.

Calculer $T=\int\int_D (x+y)^2e^{x^2-y^2}dxdy$ en utilisant le changement de variables : $x=\frac{u+v}{2}; y=\frac{u-v}{2}$.


4. Soit $\alpha> 0$ et $\beta>0$ des réels données. Soit

$D= \lbrace (x,y,z)\in \mathbb{R}^3, 0\leq x\leq 1, 0\leq y\leq 1,0\leq z\leq xy\rbrace$.

Calculer $T=\int\int\int_D x^\alpha y^\beta zdxdydz$.}

\medskip{\textbf{Exercice 2}

Soit $E$ un espace vectoriel de dimension finie $n$ et $f$ un endomorphisme de $E$.
 
 1. On suppose que $f$ est nilpotent, c'est-à-dire qu'une certaine puissance de $f$ est nulle.
 
 On appelle indice de nilpotence de $f$ le plus petit entier $p$ vérifiant : $f^p = 0$ mais
 $f^{p-1}\neq 0$.
 
(a) Montrer que si $x$ est un vecteur vérifiant $f^{p-1}(x)= 0$, alors la famille $(x,f(x),...,f^{p-1}(x))$ est une famille libre de $E$.
 
 (b) En déduire qu'on a toujours $p \leq n$, et que $f^n = 0$.
 
 (c) On suppose que l'indice de nilpotence de $f$ est $n$. Montrer qu'il existe $x_0 \in E$ tel que $(x_0,...,f^{n-1}(x_0))$ soit une base de $E$.
 
 2. On suppose que pour tout $x$ de $E$ il existe $p_x \in \mathbb{N}$ tel que $f^{p_x}(x) = 0$. Montrer que
 $f^n =0$}.

\medskip{\textbf{Exercice 3}

On joue à pile ou face avec une pièce non équilibrée. A chaque lancer, la probabilité d'obtenir pile est 2/3, et donc celle d'obtenir face est 1/3. Les lancers sont supposées indépendants, et on note $X$ la variable aléatoire réelle égale au nombre de lancers nécessaires pour obtenir, pour la première fois, deux piles consécutifs. Pour $n \geq 1$, on note $p_n$ la probabilité $P(X = n)$.

(1) Expliciter les évènements $(X = 2), (X = 3), (X = 4)$, et déterminer la valeur de $p_2,
 p_3, p_4$.

(2) Montrer que l'on a $p_n = \frac{2}{9}p_{n-2}+\frac{1}{3}p_{n-1}, n\geq 4.$
 
 (3) En déduire l'expression de $p_n$ pour tout $n$.
 
 (4) Rappeler, pour $q \in ]-1,1[$, l'expression de       
$\sum_{n=0}^{+\infty}nq^n$ et calculer alors $E(X)$}.

\medskip\textbf{Exercice 4}

Le paramètre $a$ est un réel strictement positif, $a > 0$. 

Soit une application $f : [0,a]\longrightarrow \mathbb{R}$ continue, vérifiant pour tout réel $x, 0 \leq x \leq a$, les deux conditions suivantes : 

$f(x)\neq -1$

$f(x).f(a-x)=1$

 Calculer l'intégrale $I = \int_{0}^{a}\frac{1}{1+f(x)}dx$

\quad
\quad

\end{spacing}
\end{document}
